\clearpage
\phantomsection\addcontentsline{toc}{section}{References}
\begin{thebibliography}{99}\small\raggedright
\bibitem{PaperI} UNNS Research Collective, \emph{Arithmetic Transport and Dimensional Obstructions in UNNS Occurrence Geometry}, Manuscript~I of the programme GEO-HD-1, version~0.2 (2~Nov.\ 2026); file \texttt{PAPER\_I\_TRANSPORT\_v0.2.pdf}, SHA-256 \texttt{5f7c0809db2006b419c7f7490c1aca01e984b5ad5f814ae3209e44728f312a68}.
\bibitem{PaperII} UNNS Research Collective, \emph{Generative Histories, Arithmetic Recombination, and Higher-Dimensional Route Topology}, Manuscript~II of the programme GEO-HD-1, version~0.2 (2~Nov.\ 2026); file \texttt{PAPER\_II\_HISTORIES\_v0.2.pdf}, SHA-256 \texttt{55c38d4ebd217267114381ae8d46b103a0878b940f97e6ee6cc054a7cd8cd98e}.
\bibitem{FoundI} UNNS Substrate programme, \emph{Mathematical Foundations of the UNNS Substrate, Volume~I: Generative Arithmetic and Structural Reachability} (the grammar SELECT-II and the occurrence set). Source: \texttt{ukbbi/UNNS}, commit \texttt{9a021346b2b3d964523440d2737b9d0c4ddbb3e5}, \texttt{projects/UNNS\_MATH\_FOUNDATIONS/Manuscript}.
\bibitem{FoundII} UNNS Substrate programme, \emph{Mathematical Foundations of the UNNS Substrate, Volume~II: Collective Dynamics, Prime Rigidity, and Effective Predictability} (the collective map FIELD-I and its theorems cited here, e.g.\ the eventual rigidity theorem for $G_\infty=1$). Same commit and path as~\cite{FoundI}.
\bibitem{GEOI} GEO-I, Papers~A and~B of the geometry programme (occurrence geometry, region-Graver sets and the realised complexes), cited at commit \texttt{9a021346b2b3d964523440d2737b9d0c4ddbb3e5} as in~\cite{FoundI}.
\bibitem{RefColl} Foundations reference collection v001 (SELECT-I/II, FIELD-I v001--v012; frozen), \texttt{ukbbi/UNNS}, commit \texttt{bb8b12d753c600c597feb00df17fe139c7cc65fa} (23 Oct.\ 2026), \texttt{projects/UNNS\_MATH\_FOUNDATIONS/sources}.
\bibitem{HD1} GEO-HD-1 research collection, snapshot \texttt{a88f31487f3b05c9db297dcd1192d56eea2b0c7d} (the pin named in the programme specification; not independently verified in this work). The release archives from which the records of this manuscript are taken, with their SHA-256 as computed on the files supplied for this work:
\begin{itemize}[leftmargin=1.4em,itemsep=1pt,topsep=2pt]\footnotesize
\item\texttt{GEO\_HD\_1D\_v001} \texttt{c532b7410f3bf9d437f9f172b3e56029cd9f7bab5b2ac0c0e533713155eaf7a4};
\item\texttt{GEO\_HD\_1D\_v002} \texttt{8baf463c07bb5285dd53b355aee05a0f841c068ba3ae4794f3702adff6a4bae0};
\item\texttt{GEO\_HD\_1D\_v003} \texttt{65dbb281e4c75d625b6cf3e4bcb8262d7d42e68d0c476fbba225fcde744a4706};
\item\texttt{GEO\_HD\_1D\_v004} \texttt{e0390e2a2f88eaeb11042de09fd056741ed98f94318de708d308faa090f05b5b};
\item\texttt{GEO\_HD\_BRIDGE\_v001} \texttt{e99472809928f4a23fb292f60759965159c7a4db5f018d56748ccac328719395};
\item\texttt{GEO\_HD\_BRIDGE\_v002} \texttt{012ead97e5b8854cd5218417d46f0e388d927b6d59d5f6e75a8cc343f5dd4955} (with the accepted erratum \texttt{GEO\_HD\_BRIDGE\_v002\_ERRATUM\_B2-T9.md}; the erratum for B2-T7 is \texttt{GEO\_HD\_BRIDGE\_v002\_ERRATUM\_B2-T7.md});
\item\texttt{GEO\_HD\_BRIDGE\_v003} \texttt{6fba87af9f3096772e43f4ac77fa8d55aea9e242eeef241081a987b567d5d5a8};
\item\texttt{GEO\_HD\_BRIDGE\_v004} \texttt{cfb31157a39666d1a4bbb33b280ec626c7de2c141e8232da43c8b934b8affff4};
\item\texttt{GEO\_HD\_BRIDGE\_v005} \texttt{50e342821d7a5f760c5127df348e97f3a4d4d3bf34bf60e1019d6336f28c5d03}.
\end{itemize}
\bibitem{Dirichlet1837} P.~G.~L. Dirichlet, Beweis des Satzes, dass jede unbegrenzte arithmetische Progression, deren erstes Glied und Differenz ganze Zahlen ohne gemeinschaftlichen Factor sind, unendlich viele Primzahlen enth\"alt, \emph{Abh.\ K\"onigl.\ Preuss.\ Akad.\ Wiss.} (1837), 45--81. Used for primes in a residue class $a\bmod m$ with $\gcd(a,m)=1$ (\cref{thm:R3T1}, \cref{thm:A4T1}).
\bibitem{Graver1975} J.~E. Graver, On the foundations of linear and integer linear programming~I, \emph{Math.\ Programming} \textbf{9} (1975), 207--226.
\bibitem{Sturmfels1996} B.~Sturmfels, \emph{Gr\"obner Bases and Convex Polytopes}, University Lecture Series~8, American Mathematical Society, 1996 (Graver bases, \S4).
\bibitem{Hatcher2002} A.~Hatcher, \emph{Algebraic Topology}, Cambridge University Press, 2002 (cellular homology, \S2.2; long exact sequence of a pair, \S2.1; quotient by a contractible subcomplex, Prop.~0.17; nerve theorem, Cor.~4G.3).
\bibitem{CarlssonZomorodian2009} G.~Carlsson, A.~Zomorodian, The theory of multidimensional persistence, \emph{Discrete Comput.\ Geom.} \textbf{42} (2009), 71--93.
\bibitem{CarlssonDeSilva2010} G.~Carlsson, V.~de~Silva, Zigzag persistence, \emph{Found.\ Comput.\ Math.} \textbf{10} (2010), 367--405 (zig-zag diagrams of complexes; cited for the shape $K\leftarrow K^\circ\to K'$ only).
\bibitem{Stanley1997} R.~P. Stanley, \emph{Enumerative Combinatorics}, vol.~1, Cambridge University Press, 1997 (quasi-polynomials and Ehrhart theory, \S4.4).
\bibitem{Ehrhart1962} E.~Ehrhart, Sur les poly\`edres rationnels homoth\'etiques \`a $n$ dimensions, \emph{C.~R.\ Acad.\ Sci.\ Paris} \textbf{254} (1962), 616--618.
\bibitem{Eisenbud1995} D.~Eisenbud, \emph{Commutative Algebra with a View Toward Algebraic Geometry}, Graduate Texts in Mathematics~150, Springer, 1995 (Hilbert functions of finitely generated graded modules, Chapter~12; Noetherianity).
\bibitem{Kozlov2008} D.~N. Kozlov, \emph{Combinatorial Algebraic Topology}, Algorithms and Computation in Mathematics~21, Springer, 2008.
\end{thebibliography}
\paragraph{On the references.} UNNS sources are cited at the commits and archive hashes fixed by the programme; repository access during this work was read-only and the pinned repositories were not fetched from the environment in which the manuscript was prepared. External literature is cited by author, title and publication data and only for the standard theorems named (Dirichlet's theorem, cellular and relative homology, the long exact sequence, the nerve theorem, Hilbert functions and quasi-polynomials) and for the resemblance of the zig-zag shape and of Graver bases to known notions; nothing in the manuscript is proved by appeal to an equivalence with those theories, and no claim of novelty or priority is made.
